\documentclass{article}
\usepackage[T1]{fontenc}
\usepackage{lmodern}
\usepackage{graphicx}
\usepackage{amsmath,amssymb}
\usepackage[a4paper,margin=2cm]{geometry}
\usepackage[absolute,overlay]{textpos}
\setlength{\TPHorizModule}{1mm}
\setlength{\TPVertModule}{1mm}
\usepackage[indonesian]{babel}
\usepackage{hyperref}
\setlength{\emergencystretch}{2em}
% unit: Halaman 36

\begin{document}

% === Halaman 36 (llm) ===

\section*{2. Kode Program RSA Python (2)}

\begin{verbatim}
from Crypto.PublicKey import RSA
from Crypto.Cipher import PKCS1_OAEP

def PembangkitanKunciRSA():
    # Buat kunci RSA
    key = RSA.generate(1024)
    kunci_publik = key.pubkey()

    # Tampilkan informasi kunci
    print("Modulus (n):", key.n)
    print("Kunci publik (e):", key.e)
    print("Kunci privat (d):", key.d)
    print("Bilangan prima p:", key.p)
    print("Bilangan prima q:", key.q)

    # Simpan kunci publik dan kunci privat ke dalam dua file kunci berbeda dengan format PEM
    with open("kunci_publik.pem", "wb") as file:
        file.write(kunci_publik.exportKey('PEM'))
        file.close()

    with open("kunci_privat.pem", "wb") as file:
        file.write(key.exportKey('PEM','passwordku'))  # Password bisa diubah
        file.close()
\end{verbatim}

\end{document}
